\section{Heterogeneidad y efectos distributivos}\label{sec:heterogeneity}

\subsection{¿Quién se ajusta?}

La Figura~\ref{fig:het} reporta la estimación de diferencias en diferencias entre
trabajadores poco calificados y calificados para cada resultado principal dentro de
subgrupos demográficos y regionales.\footnote{Cada estimación corresponde a las diferencias
en diferencias dentro del subgrupo, de modo que pregunta si el contraste entre poco
calificados y calificados difiere, por ejemplo, entre hombres y mujeres, y no cuál es el
efecto en niveles para un subgrupo.} Destacan tres patrones, coherentes con la intuición
económica sobre dónde el salario mínimo es más vinculante.

Primero, ningún subgrupo muestra ganancias salariales. Ya sea que se definan por sexo,
edad, área o incidencia (\emph{bite}) regional, todas las estimaciones salariales por
subgrupo son negativas o indistinguibles de cero. La ausencia de respuesta en las
remuneraciones es un rasgo generalizado de los datos y no un promedio que esconda efectos
que se compensan.

Segundo, el ajuste de la composición crece con la exposición. La caída del empleo formal
es mayor en los departamentos de alta incidencia ($-5.5$ puntos porcentuales) que en los de
baja incidencia ($-3.6$ puntos), y lo mismo ocurre con el aumento del trabajo independiente
($+5.0$ frente a $+2.8$ puntos). Este patrón de dosis y respuesta es el que debería
producir un efecto causal del piso salarial, y es la contraparte transversal de los
resultados de exposición continua de la Sección~\ref{sec:robustness}. La estimación del
nivel de empleo, en cambio, \emph{no} crece con la incidencia (es menor y no significativa
en los departamentos de alta incidencia), lo que constituye una razón adicional para
atribuirla a la recuperación diferencial y no a la reforma.

Tercero, la reasignación se concentra entre los jóvenes y las mujeres. La caída del empleo
formal es mayor para los trabajadores de catorce a veinticinco años ($-5.8$ puntos) y para
las mujeres ($-5.5$ puntos), los grupos cuyos salarios tienen mayor probabilidad de
ubicarse cerca del piso y cuyo vínculo con el empleo formal es más frágil; el aumento del
trabajo independiente es el doble de grande para las mujeres que para los hombres. Para los
trabajadores mayores de cuarenta y cinco años, el efecto sobre la formalidad es cercano a
cero, aunque su trabajo independiente también aumenta. Esto replica el énfasis en los
jóvenes de la literatura peruana previa \citep{cespedes_2005, jaramillo_lopez_2006} y el
énfasis en las mujeres y en los menos calificados de la evidencia para países en
desarrollo \citep{katzkowicz_2021_uruguay, lemos_2009}.

\begin{figure}[!t]
\centering
\includegraphics[width=0.98\textwidth]{fig6_heterogeneity.pdf}
\caption{Efectos heterogéneos de la reforma de 2022. Cada punto es la estimación de
diferencias en diferencias $\mathrm{Low}\times\mathrm{Post}$ dentro del subgrupo indicado,
con intervalos de confianza al noventa y cinco por ciento basados en errores estándar
agrupados por departamento.}
\label{fig:het}
\end{figure}

\subsection{Márgenes de ajuste}

¿Dónde ocurre exactamente la reasignación? Cuatro márgenes orientados a los mecanismos
precisan el panorama (Cuadro~\ref{tab:margins} del Apéndice). Primero, la proporción de
asalariados de baja educación empleados \emph{sin contrato escrito} aumenta en 3.1 puntos
porcentuales respecto de los asalariados de alta educación, de modo que la informalización
opera dentro del empleo asalariado y no solo mediante salidas hacia el trabajo
independiente.\footnote{La pregunta sobre el contrato se eliminó del cuestionario reducido
que se aplicó en quince regiones en febrero de 2021; codifico esas observaciones como
faltantes y no como carentes de contrato, lo que importa para este margen.} Segundo, el
empleo se recompone hacia las \emph{microunidades}: la probabilidad de que un trabajador de
baja educación esté empleado en un establecimiento de a lo sumo veinte trabajadores
aumenta en 2.8 puntos, y las unidades pequeñas son precisamente donde la fiscalización es
más escasa. Tercero, según la descomposición del empleo informal del INEI (disponible hasta
2023), el margen desplazado recae en el \emph{sector informal} propiamente dicho (+2.4
puntos) y no en empleos no registrados dentro de empresas formales (+0.8, no
significativo): los trabajadores no se quedaron con empleadores formales de manera
irregular, sino que se trasladaron a unidades productivas informales. Cuarto, la
proporción de asalariados con menos de un año de antigüedad no se mueve, de modo que el
ajuste no se manifiesta como un congelamiento de contrataciones entre quienes siguen siendo
asalariados.

Dos divisiones adicionales por subgrupo ubican la carga. Los efectos de composición se
concentran casi por completo entre los \emph{no jefes} de hogar (perceptores secundarios de
ingreso, cuya formalidad cae 6.8 puntos, frente a un efecto nulo para los jefes de hogar) y
entre los trabajadores no indígenas, lo que es coherente con los mercados laborales urbanos
de la costa, donde se concentra el empleo formal en torno al piso. En suma, los empleos
formales marginales que se disolvieron fueron los de perceptores secundarios en pequeñas
empresas urbanas.

\subsection{Efectos salariales distributivos}

La ausencia de un efecto salarial promedio podría, en principio, ocultar movimientos que se
compensan en distintos puntos de la distribución salarial, dado que se espera que un
salario mínimo comprima la cola inferior. Para examinarlo, estimo regresiones por cuantiles
incondicionales (función de influencia recentrada) a la manera de
\citet{firpo_fortin_lemieux_2009}, aplicando la especificación de diferencias en
diferencias a la función de influencia de cada decil del logaritmo del salario real por
hora. La Figura~\ref{fig:rif} grafica el coeficiente $\mathrm{Low}\times\mathrm{Post}$ a lo
largo de los deciles. Las estimaciones son no positivas en todos los deciles, en su mayoría
entre $-2$ y $-4$ por ciento y significativas en algunos deciles intermedios; lo crucial es
que no hay ningún efecto positivo en los deciles inferiores, donde un salario mínimo
vinculante y fiscalizado comprimiría la distribución desde abajo. La evidencia distributiva
refuerza, por lo tanto, la ausencia de respuesta en las remuneraciones: la reforma no elevó
la cola inferior de la distribución salarial de los poco calificados respecto de la de los
calificados. Esto es coherente con la distribución de ingresos de la Figura~\ref{fig:dist},
que muestra poco movimiento hacia el nuevo piso, y con la interpretación de que, en un
mercado laboral con incumplimiento generalizado, el mínimo legal no fue vinculante para la
mayoría de los trabajadores poco calificados.

\begin{figure}[!t]
\centering
\includegraphics[width=0.72\textwidth]{fig8_rif_quantile.pdf}
\caption{Efectos salariales distributivos. Cada punto es el coeficiente de diferencias en
diferencias $\mathrm{Low}\times\mathrm{Post}$ de una regresión por cuantiles
incondicionales (RIF) en el decil indicado del logaritmo del salario real por hora
\citep{firpo_fortin_lemieux_2009}. Las bandas son intervalos de confianza al noventa y
cinco por ciento con errores agrupados por departamento.}
\label{fig:rif}
\end{figure}
